\section{Discussion}
\label{sec:discussion}

This section interprets the results against the five hypotheses and the seven objectives; the numbers are in
Section~\ref{sec:results}.

\textbf{H1, tiles against whole-photo input: supported on the headline split.} With the same encoder and head, tiles at
a common leaf scale beat one whole-photo description at both sizes. Two further observations point the same way:
whole-image scores rise with input size (Table~\ref{tab:final_conditions}), and tiles at the stored scale do worse than
tiles at the common scale (Table~\ref{tab:final_ablations-test}). Detail matters. It does not follow that tiling matters:
ConvNeXt-T at 448 px has the higher mean and the paired test finds no clear difference. What the tile pipeline offers is a
different cost profile, a small head trained from cached descriptions, against fine-tuning a whole network at
a larger input size. The Committee gain suggests that the two kinds of model make partly different errors.
On the locked split the benefit over the 512-px control was not clear, which is a reminder that the split changes
conclusions.

\textbf{H2, attention against mean, max and Otsu: supported.} Most tiles hold background, and the head gives them far less than their share of its attention (Section~\ref{sec:results}). Mean pooling averages all of them with the few informative
tiles, which is a likely reason for its low score; max pooling keeps the single strongest tile and comes close, without a clear
difference on balanced accuracy. Hard background removal by Otsu's method did not help whole-image models and lowered the
tile model, so a learned weighting looked more useful than a fixed mask; an earlier report found Otsu alone changed
accuracy little \cite{ali2026}.

\textbf{H3, frozen features on field photos: not shown.} All methods score low on Paddy Doctor and none is clearly
ahead of the fine-tuned models, so in-domain scores say little about field behaviour. Matching the zoom did not
help. We did not test the likely causes, such as lighting, camera and field context, and the primary photos have plain
backgrounds (Section~\ref{sec:setup:data}); we therefore do \emph{not} claim robustness to clutter.

\textbf{H4, heat maps point at symptoms: supported as a qualitative guide.} The deletion test shows that the decision
rests on the attended tiles, but that alone does not show that those tiles are symptoms. The ratings are mixed, and about a
quarter of the ratings in each set say no. The effect on field photos is small. We treat heat maps as an aid for a reader,
not as a segmentation, and make no severity claim.

\textbf{H5, cost: a descriptive answer.} The pipeline is slower per photo than whole-image models because of its many
encoder passes. Fewer tiles and a distilled head reduce the cost at little loss of validation score, but the 8-bit encoder
fails, so an on-device build is not shown.

\textbf{What the shortcut and split checks mean for reading in-domain numbers.} The nearest-in-time baseline matches
the models on the locked split, and the partition-order experiment shows that the order of augmenting and splitting
moves a score more than most differences between methods. We therefore read differences between methods only on the
session-blocked split and only through the paired rule, and we state the split beside every number. This is a general
lesson for the field rather than a verdict on any study: the details of splitting are rarely reported
\cite{xu2024}. The per-class results add that hispa and leaf blast are hard for every method, so a single accuracy figure
hides the classes where errors concentrate; hispa photos are most often called healthy.

\textbf{Objectives.} Table~\ref{tab:objectives} relates the seven objectives to the evidence.

\begin{table}[t]
\caption{The seven objectives, the evidence and the status.}
\label{tab:objectives}
\centering
\scriptsize
\renewcommand{\arraystretch}{1.1}
\begin{tabular}{@{}L{3.1cm}L{3.5cm}L{1.1cm}@{}}
\toprule
Objective & Evidence & Status \\
\midrule
1 Two-dataset benchmark & Primary photos and Paddy Doctor (Section~\ref{sec:setup:data}) & met \\
2 Leakage-free protocol & Grouped splits made first, mean $\pm$ SD, two ranges, one rule (Section~\ref{sec:setup}) & met \\
3 Whole-image baselines, including Otsu & Four models at 128, 224 and 448 px, with and without Otsu (Tables~\ref{tab:final_main}, \ref{tab:final_conditions}) & met \\
4 Tile-based weakly supervised model & Attention MIL (Section~\ref{sec:method}) & met \\
5 Ablations of resolution, aggregation, Otsu & Tables~\ref{tab:final_conditions}, \ref{tab:final_main}, \ref{tab:final_ablations-test}, \ref{tab:final_ablations} & met \\
6 Zero-shot transfer; attention against symptoms & Paddy Doctor (H3 not shown); deletion test and model-rater ratings; no expert or outline check & partly met \\
7 Cost and compressibility & Table~\ref{tab:final_cost}; fewer tiles and distillation help, 8-bit fails; no on-device build & partly met \\
\bottomrule
\end{tabular}
\end{table}
